% GENERATED LOCAL-BUILD FLAT PART — DO NOT MERGE

% ---- BEGIN TIR/monograph/v12/chapters/ch08_three_flavour_carrier_su3.tex ----
% !TEX root = ../../tir_monograph_v12.tex
% V12_SOURCE_MAP:
% - TIR/foundations/TIR_KAPPA_FLAVOUR_MIXING_NORMALIZATION_V0_1.md
% - TIR/monograph/appendices/appC_su3_primer.tex
% - TIR/monograph/chapters/ch16_pmns_mixing.tex
% - TIR/monograph/chapters/ch18_ckm_matrix.tex
\chapter[Three-Flavour Carrier]{Three-Flavour Carrier and $SU(3)_F$}
\label{ch:v12_three_flavour_su3}

\paragraph{Established context.}
For a complex three-component carrier, unitary basis changes are described by
$U(3)$, while determinant-one unitary transformations form $SU(3)$.  The Lie
algebra dimension
\[
\dim_{\mathbb R}\mathfrak{su}(N)=N^2-1
\]
is standard.  TIR uses this group-theoretic structure as the family-mixing
carrier for its three-flavour branch.

\vTwelveStatus{B}{--}{PASS}

\section{Three-flavour carrier}
The family branch is represented by
\begin{equation}
\boxed{V_F\cong\mathbb C^3.}
\label{eq:v12-flavour-carrier}
\end{equation}
Choose an ordered flavour basis
\[
\{|f_1\rangle,|f_2\rangle,|f_3\rangle\}.
\]
A normalized family state is
\[
|\psi_F\rangle
=\sum_{i=1}^{3}c_i|f_i\rangle,
\qquad
\sum_{i=1}^{3}|c_i|^2=1.
\]
Full determinant-one family mixing acts through
\begin{equation}
\boxed{U_F\in SU(3)_F.}
\label{eq:v12-su3f-carrier}
\end{equation}
The carrier dimension therefore fixes the flavour multiplicity
\begin{equation}
\boxed{N_F=3.}
\label{eq:v12-flavour-multiplicity}
\end{equation}

\section{Eight independent mixing directions}
The associated mixing algebra is
\[
\mathfrak{su}(3)_F,
\]
with
\begin{equation}
\boxed{
\dim_{\mathbb R}\mathfrak{su}(3)_F
=3^2-1
=8.
}
\label{eq:v12-su3f-dimension}
\end{equation}
If $\{T_a\}_{a=1}^{8}$ is any basis of generators, a local family
transformation may be written
\[
U_F(\boldsymbol\theta)
=\exp\!\left(i\sum_{a=1}^{8}\theta_aT_a\right).
\]
The number eight is basis-independent: it is the dimension of the Lie algebra,
not a property of a particular generator convention.

The TIR symmetric-pair crosswalk supplies the decomposition
\begin{equation}
\mathfrak{su}(3)_F
=\mathfrak{so}(3)\oplus\mathfrak p,
\qquad
8=3+5,
\label{eq:v12-su3-symmetric-pair}
\end{equation}
which gives an independent carrier-level accounting of the same eight real
mixing directions.

\section{Generator--flavour incidence set}
For later normalization bookkeeping define
\begin{equation}
\mathcal C_{\rm mix}
=\{1,\ldots,8\}_{\rm generator}
\times\{1,2,3\}_{\rm flavour}.
\label{eq:v12-mixing-incidence-set}
\end{equation}
Its cardinality is
\begin{equation}
|\mathcal C_{\rm mix}|
=8\cdot3
=24.
\label{eq:v12-mixing-incidence-cardinality}
\end{equation}
At this stage the integer $24$ is only the exact incidence count generated by
the already specified carrier and algebra.  Chapter~\ref{ch:v12_kappa_normalization}
adds the half-turn phase measure and the binary-information numerator to obtain
the TIR normalization.

\section{Downstream mixing matrices}
The PMNS and CKM constructions are treated in
Chapter~\ref{ch:v12_ckm_pmns}.  Their numerical angles, phases and empirical
residuals are downstream sector assignments.  The carrier statements
\eqref{eq:v12-flavour-carrier}--\eqref{eq:v12-su3f-dimension} are the structural
parents shared by those sectors and by the normalization chapter that follows.

% ---- END TIR/monograph/v12/chapters/ch08_three_flavour_carrier_su3.tex ----
\clearpage

% ---- BEGIN TIR/monograph/v12/chapters/ch09_kappa_normalization.tex ----
% !TEX root = ../../tir_monograph_v12.tex
% V12_CANONICAL_KAPPA_OWNER
% Canonical theorem source:
% TIR/foundations/TIR_KAPPA_FLAVOUR_MIXING_NORMALIZATION_V0_1.md
% Validator:
% TIR/validation/tir_kappa_flavour_mixing_normalization_v0_1.py
\chapter[The Kappa Normalization]{The \texorpdfstring{$\kappa=\ln2/(24\pi)$}{kappa = ln2/(24 pi)} Normalization}
\label{ch:v12_kappa_normalization}

\paragraph{Established context.}
The binary-information identity $H_2(1/2)=\ln2$ is standard Shannon
information theory \cite{shannon1948,CoverThomas2006}.  The angular closure of
one full turn is $2\pi$, and $\dim\mathfrak{su}(3)=8$ is the standard
special-unitary Lie-algebra dimension.  TIR combines these established
ingredients with the three-flavour carrier derived in the preceding chapter.

\vTwelveStatus{B}{--}{PASS}

\section{Binary information numerator}
For a binary distribution $(1-p,p)$,
\[
H_2(p)=-p\ln p-(1-p)\ln(1-p).
\]
At the balanced relational coordinate,
\begin{equation}
\boxed{
\sigma_\star=\frac12,
\qquad
I_\star=H_2(1/2)=\ln2.
}
\label{eq:v12-kappa-binary-information}
\end{equation}
This supplies the informational numerator.

\section{Mixing-channel multiplicity}
From Chapter~\ref{ch:v12_three_flavour_su3},
\[
N_F=3,
\qquad
\dim_{\mathbb R}\mathfrak{su}(3)_F=8.
\]
The generator--flavour incidence set therefore has
\begin{equation}
\boxed{
N_{\rm mix}
=N_F\dim\mathfrak{su}(3)_F
=3\cdot8
=24.
}
\label{eq:v12-kappa-mixing-count}
\end{equation}
Equivalently,
\[
N_{\rm mix}
=N_F(N_F^2-1)
=3(3^2-1)
=24.
\]

\section{Half-turn phase unit}
The primitive half coordinate acts on the standard angular closure as
\begin{equation}
\boxed{
\Delta\phi_{1/2}
=\frac12(2\pi)
=\pi.
}
\label{eq:v12-kappa-half-turn}
\end{equation}
Assigning one such phase unit to each primitive generator--flavour incidence
channel gives the total mixing-phase measure
\begin{equation}
\boxed{
\Phi_{\rm mix}
=N_{\rm mix}\Delta\phi_{1/2}
=24\pi.
}
\label{eq:v12-kappa-mixing-phase-measure}
\end{equation}
Thus
\begin{equation}
\boxed{
24\pi
=\underbrace{(3^2-1)}_{8\ \mathrm{mixing\ directions}}
\underbrace{3}_{\mathrm{flavours}}
\underbrace{\pi}_{\mathrm{half\!\!-turn\ phase}}.
}
\label{eq:v12-kappa-denominator-factorization}
\end{equation}

\section{Structural normalization}
The TIR information-per-mixing-phase normalization is
\begin{equation}
\boxed{
\kappa
=\frac{I_\star}{\Phi_{\rm mix}}
=\frac{\ln2}{N_F(N_F^2-1)\pi}
=\frac{\ln2}{3(3^2-1)\pi}
=\frac{\ln2}{24\pi}.
}
\label{eq:v12-kappa-canonical}
\end{equation}
Numerically,
\[
\kappa
=0.009193150006360484\ldots .
\]
Equation~\eqref{eq:v12-kappa-canonical} is the canonical v12 publication
location of the full derivation.  Downstream chapters refer to this equation and
to its theorem/validator provenance.

\section{Independent tetrahedral integer crosscheck}
The spatial branch obtains the regular tetrahedral finite cell with abstract
automorphism group
\[
\operatorname{Aut}(\Delta^3)\cong S_4,
\qquad
|S_4|=24.
\]
Hence
\begin{equation}
\boxed{
3\,\dim\mathfrak{su}(3)_F
=3\cdot8
=24
=|S_4|.
}
\label{eq:v12-kappa-tetra-crosscheck}
\end{equation}
The flavour-mixing incidence route supplies the normalization dependency; the
spatial route supplies an independent finite-symmetry occurrence of the same
integer.

\section{Exact phase-rate consequence}
Let
\[
\omega\equiv\frac{d\phi}{dt}=2\pi f,
\qquad
d\mathcal I\equiv\kappa\,d\phi.
\]
Then
\begin{equation}
\boxed{
\Gamma_{\mathcal I}
\equiv\frac{d\mathcal I}{dt}
=\kappa\omega
=\frac{\ln2}{12}f.
}
\label{eq:v12-kappa-phase-rate}
\end{equation}
One complete angular cycle carries
\begin{equation}
\boxed{
\Delta\mathcal I_{\rm cycle}
=2\pi\kappa
=\frac{\ln2}{12}.
}
\label{eq:v12-kappa-information-cycle}
\end{equation}

\section{Constraint manifold}
For
\[
\mathbf q=(\kappa,\omega,f,\Gamma_{\mathcal I})
\]
and constraints
\[
C_1=\kappa-\frac{\ln2}{24\pi}=0,
\qquad
C_2=\omega-2\pi f=0,
\qquad
C_3=\Gamma_{\mathcal I}-\kappa\omega=0,
\]
the constraint Jacobian has rank three.  The declared subsystem is therefore
one-dimensional and can be parameterized by
\begin{equation}
\boxed{
\mathbf q(f)=
\left(
\frac{\ln2}{24\pi},
2\pi f,
f,
\frac{\ln2}{12}f
\right).
}
\label{eq:v12-kappa-constraint-curve}
\end{equation}

\section{Reproducibility}
The theorem source is
\path{TIR/foundations/TIR_KAPPA_FLAVOUR_MIXING_NORMALIZATION_V0_1.md}.
Its deterministic numerical/algebraic validator is
\path{TIR/validation/tir_kappa_flavour_mixing_normalization_v0_1.py}.

% ---- END TIR/monograph/v12/chapters/ch09_kappa_normalization.tex ----
\clearpage

% ---- BEGIN TIR/monograph/v12/chapters/ch10_discrete_structural_labels.tex ----
% !TEX root = ../../tir_monograph_v12.tex
% V12_SOURCE_MAP:
% - TIR/monograph/chapters/ch03_l_constants.tex
% - TIR/monograph/chapters/ch04_quark_primes.tex
% - TIR/STRUCTURAL_CHOICES.md
% - TIR/foundations/TIR_PLATONIC_L_CONSTANTS_CLOSURE_V0_1.md
% Validators:
% - TIR/validation/tir_v12_discrete_labels_audit_v0_1.py
% - TIR/validation/tir_platonic_l_constants_closure_v0_1.py
\chapter[Discrete Structural Labels]{Discrete Structural Labels -- $L_i$, Primes and Provenance}
\label{ch:v12_discrete_labels}

\paragraph{Established context.}
The Collatz map, prime-number sequence, Platonic classification and finite groups
\(A_4,S_4,A_5\) are standard mathematical objects.  TIR uses selected finite
properties of those objects inside an explicitly typed structural layer.  The
v12.1 synchronization separates three levels: standard mathematical facts, the
TIR selection/closure rule, and any downstream physical interpretation.

\section{Platonic index alphabet}

For a convex regular polyhedron with Schlaefli symbol \(\{p,q\}\),
\[
\frac1p+\frac1q>\frac12,
\qquad p,q\ge3.
\]
The integer solutions are exactly
\[
\boxed{
\{3,3\},\{3,4\},\{3,5\},\{4,3\},\{5,3\}.
}
\]
Thus the complete index alphabet appearing in the convex Platonic
classification is
\[
\boxed{\{3,4,5\}}.
\]
The triangular-face branch is
\[
\{3,3\},\qquad\{3,4\},\qquad\{3,5\},
\]
with orientation-preserving rotational groups
\[
G_3\simeq A_4,
\qquad
G_4\simeq S_4,
\qquad
G_5\simeq A_5,
\]
and orders
\[
|A_4|=12,
\qquad |S_4|=24,
\qquad |A_5|=60.
\]
Equivalently, the spherical triangle groups \((2,3,q)\) have
\[
|G_q|
=\frac{2}{\frac12+\frac13+\frac1q-1},
\qquad q\in\{3,4,5\}.
\]
At \(q=6\), the spherical excess vanishes; hence the finite spherical branch
terminates at \(q=5\).

\section{Common tetrahedral root and subgroup indices}

The tetrahedral rotational group satisfies
\[
A_4\triangleleft S_4,
\]
so
\[
\boxed{[S_4:A_4]=\frac{24}{12}=2}.
\]
A point stabilizer in the natural five-point action of \(A_5\) is isomorphic to
\(A_4\), giving
\[
\boxed{[A_5:A_4]=\frac{60}{12}=5}.
\]
Define the two extension coset spaces over the tetrahedral root,
\[
X_4:=S_4/A_4,
\qquad
X_5:=A_5/A_4.
\]
Then
\[
|X_4|=2,
\qquad
|X_5|=5.
\]
TIR identifies
\[
\boxed{L_4:=|X_4|=2},
\qquad
\boxed{L_5:=|X_5|=5}.
\]

\section[TIR root-extension closure and L3]{TIR root-extension closure and $L_3$}

The TIR structural step is explicit:

\begin{quote}
For the tetrahedral root, the complete non-self-dual Platonic extension carrier
is the disjoint union of all admissible extension coset spaces.
\end{quote}

Therefore
\[
X_3^{\rm closure}:=X_4\sqcup X_5
=S_4/A_4\sqcup A_5/A_4.
\]
Since the union is disjoint,
\[
|X_3^{\rm closure}|=|X_4|+|X_5|=2+5=7,
\]
and TIR defines
\[
\boxed{L_3:=|X_3^{\rm closure}|=7}.
\]
Thus
\[
\boxed{
(L_3,L_4,L_5)
=\left(
|S_4/A_4|+|A_5/A_4|,
|S_4/A_4|,
|A_5/A_4|
\right)
=(7,2,5).
}
\]
\vTwelveStatus{B}{--}{PASS}

The standard finite-group facts are not fitted.  The root-extension closure
rule is a TIR structural definition; the triple follows exactly once that rule
is admitted.  The finite-group derivation alone does not establish a physical
law, a dynamical role for the labels, or a derivation from a specific Kaehler
manifold.

Canonical theorem:
\path{TIR/foundations/TIR_PLATONIC_L_CONSTANTS_CLOSURE_V0_1.md}.

Deterministic validator:
\path{TIR/validation/tir_platonic_l_constants_closure_v0_1.py}.

\section[Independent finite Collatz crosscheck for L3]{Independent finite Collatz crosscheck for $L_3$}

Define the usual Collatz step on positive integers by
\begin{equation}
C(n)=
\begin{cases}
3n+1, & n\ \text{odd},\\
n/2, & n\ \text{even}.
\end{cases}
\label{eq:v12-collatz-map}
\end{equation}
For the single seed required here, direct iteration gives the finite orbit
\begin{equation}
\boxed{
3\to10\to5\to16\to8\to4\to2\to1.
}
\label{eq:v12-collatz-three-orbit}
\end{equation}
There are exactly seven steps before the first occurrence of \(1\).  Hence the
pre-existing arithmetic provenance gives independently
\begin{equation}
\boxed{\operatorname{depth}_{\rm Collatz}(3)=7=L_3.}
\label{eq:v12-L3}
\end{equation}
\vTwelveStatus{B}{--}{PASS}
The validator evaluates only this displayed finite orbit; no global assertion
about all positive Collatz seeds enters the calculation.

\section[Independent twin-prime crosschecks for L4 and L5]{Independent twin-prime crosschecks for $L_4$ and $L_5$}

From the prime pair \((3,5)\), the older arithmetic route gives
\begin{equation}
\boxed{5-3=2=L_4,}
\qquad
\boxed{5=L_5.}
\label{eq:v12-L4-L5}
\end{equation}
\vTwelveStatus{B}{--}{PASS}
These identities are retained as independent crosschecks of the values obtained
from the Platonic subgroup-index construction, not as inputs used to fit that
construction.

Useful downstream combinations include
\[
L_3+L_4=9,
\qquad
L_3-1=6,
\qquad
\frac{L_4}{L_5}=\frac25,
\qquad
L_4^{L_5}=2^5=32.
\]
Additional exact checks connecting the Platonic index alphabet and the three
rotation-group orders are
\[
3^2+4^2=5^2,
\qquad
3+4+5=12=|A_4|,
\]
\[
2(3+4+5)=24=|S_4|,
\qquad
3\cdot4\cdot5=60=|A_5|.
\]
These are consistency checks, not extra derivation premises.

\section{Prime candidate set for quark labels}

The stronger derivation of the \(L_i\) values does not automatically derive the
quark-prime map.  That map retains its own provenance and identifiability gate.
Reserve the first prime \(p_1=2\) for the binary structural branch and define the
ordered candidate list
\begin{equation}
\mathcal P_Q
=(p_2,p_3,p_4,p_5,p_6,p_7)
=(3,5,7,11,13,17).
\label{eq:v12-prime-candidate-list}
\end{equation}
Let the ordered quark-slot list be
\begin{equation}
\mathcal F_Q=(u,d,s,c,b,t).
\label{eq:v12-quark-slot-list}
\end{equation}

\section{Arithmetic constraints and residual ambiguity}

The historical assignment satisfies
\begin{equation}
(u,d,s,c,b,t)=(3,5,7,11,13,17),
\label{eq:v12-quark-prime-assignment}
\end{equation}
together with
\[
d-u=2,
\qquad
s=L_3=7,
\qquad
|s-c|=4,
\qquad
|b-t|=4.
\]
An exhaustive scan over all
\[
6!=720
\]
bijections between \(\mathcal F_Q\) and \(\mathcal P_Q\) shows that these
arithmetic conditions alone leave exactly two survivors:
\begin{align}
(u,d,s,c,b,t)&=(3,5,7,11,13,17),\\
(u,d,s,c,b,t)&=(3,5,7,11,17,13).
\end{align}
Thus the displayed arithmetic relations leave one residual heavy-pair exchange
\(b\leftrightarrow t\).
\vTwelveStatus{D}{--}{PASS}

\section{Typed ordering rule}

The v12 structural selection is the explicit order-preserving map
\begin{equation}
\boxed{
q:\mathcal F_Q\to\mathcal P_Q,
\qquad
q((\mathcal F_Q)_i)=(\mathcal P_Q)_i,
\quad i=1,\ldots,6.
}
\label{eq:v12-quark-prime-typed-rule}
\end{equation}
Equivalently,
\[
q(u)=3,\quad q(d)=5,\quad q(s)=7,\quad q(c)=11,\quad q(b)=13,\quad q(t)=17.
\]
Within the declared ordered candidate set and ordered flavour slots, this rule
has exactly one assignment.  The uniqueness is uniqueness of the typed
selection rule; a scheme- and scale-defined physical quark-mass map remains a
separate open gate.
\vTwelveStatus{B}{--}{PASS}

\section{Reproducibility receipt}

The deterministic audit
\path{TIR/validation/tir_v12_discrete_labels_audit_v0_1.py}
checks the finite Collatz orbit and enumerates all \(720\) prime assignments.
The independent validator
\path{TIR/validation/tir_platonic_l_constants_closure_v0_1.py}
enumerates the five Platonic Schlaefli pairs, constructs the finite permutation
groups and cosets, verifies the subgroup indices and reproduces
\((L_3,L_4,L_5)=(7,2,5)\) without using the older arithmetic values as inputs.

% ---- END TIR/monograph/v12/chapters/ch10_discrete_structural_labels.tex ----
\clearpage

% ---- BEGIN TIR/monograph/v12/chapters/ch11_action_architecture_coefficient_forcing.tex ----
% !TEX root = ../../tir_monograph_v12.tex
% V12_SOURCE_MAP:
% - TIR/monograph/chapters/ch05_electron_action.tex
% - TIR/monograph/chapters/ch06_generation_release.tex
% - TIR/foundations/TIR_COEFFICIENT_ROLE_ORIENTATION_FORCING_V0_1.md
% - TIR/foundations/TIR_COEFFICIENT_MAGNITUDE_PARENT_EVALUATION_V0_1.md
% - TIR/TIR_COMPLETION_FRONTIER_V0_7.md
% Validators:
% - TIR/validation/tir_coefficient_role_orientation_forcing_v0_1.py
% - TIR/validation/tir_coefficient_magnitude_parent_evaluation_v0_1.py
\chapter[Action Architecture and Coefficient Forcing]{Action Architecture and the Coefficient-Forcing Problem}
\label{ch:v12_action_coefficients}

\section{Intrinsic coefficient generator}
The current TIR coefficient architecture is organized by the integer state
\begin{equation}
\boxed{(h,a,b,c)\in\mathbb Z^4}
\label{eq:v12-coefficient-state}
\end{equation}
and generator
\begin{equation}
\boxed{
G(h,a,b,c)
=\frac h2
+a\kappa
+b\frac{\kappa}{L_3}
+c\frac{\kappa^2}{2}.
}
\label{eq:v12-coefficient-generator}
\end{equation}
The normalization $\kappa$ is the canonical quantity derived in
Chapter~\ref{ch:v12_kappa_normalization}; the discrete label $L_3$ is fixed in
Chapter~\ref{ch:v12_discrete_labels}.

\section{Typed role routing}
Four source roles are distinguished:
\[
\mathcal R
=\{R_h,R_a,R_b,R_c\},
\]
with
\begin{align}
R_h&=\text{projective-half-spin},\\
R_a&=\text{generation-release},\\
R_b&=\text{return-axis},\\
R_c&=\text{curvature-holonomy}.
\end{align}
The corresponding generator basis is
\begin{equation}
\mathcal B
=\left\{
\frac12,
\kappa,
\frac{\kappa}{L_3},
\frac{\kappa^2}{2}
\right\}.
\label{eq:v12-generator-basis}
\end{equation}
The canonical role router gives
\begin{equation}
\boxed{
R_h\leftrightarrow h,
\qquad
R_a\leftrightarrow a,
\qquad
R_b\leftrightarrow b,
\qquad
R_c\leftrightarrow c.
}
\label{eq:v12-role-slot-map}
\end{equation}
The four parent signatures are pairwise distinct. Exhaustive permutation of the
four coefficient slots therefore leaves exactly one role-preserving
permutation: the identity.

\vTwelveStatus{B}{--}{PASS}
The deterministic validator is
\path{TIR/validation/tir_coefficient_role_orientation_forcing_v0_1.py}.

\section{Gradient orientation}
For the runtime source operator
\begin{equation}
s(\alpha_s,x)=\tanh(-\alpha_s x),
\qquad
x=\frac{\partial V}{\partial\phi},
\qquad
\alpha_s>0,
\label{eq:v12-gradient-source}
\end{equation}
monotonicity and oddness of $\tanh$ imply
\begin{equation}
\boxed{
\operatorname{sgn}s(\alpha_s,x)
=\operatorname{sgn}(-x)
}
\label{eq:v12-gradient-sign}
\end{equation}
for every finite $x$ and positive $\alpha_s$. The orientation class can
therefore be represented by
\begin{equation}
\boxed{
\chi_{\nabla}
:=\operatorname{sgn}\!\left(-\frac{\partial V}{\partial\phi}\right)
\in\{-1,0,+1\}.
}
\label{eq:v12-gradient-orientation}
\end{equation}

\vTwelveStatus{B}{--}{PASS}

\section{Directed-orbit and chiral orientation}
The directed-orbit source is
\begin{equation}
\chi_O=
\begin{cases}
+1,&\text{attractor}\to\text{satellite},\\
-1,&\text{satellite}\to\text{attractor},\\
0,&\text{neutral},
\end{cases}
\label{eq:v12-orbit-orientation}
\end{equation}
while the chiral-path source is
\begin{equation}
\chi_H=
\begin{cases}
+1,&H_+,\\
-1,&H_-,\\
0,&\text{neutral}.
\end{cases}
\label{eq:v12-chiral-orientation}
\end{equation}
For an atomic/orbital state define the active source-sign set
\[
S=\{\chi_\nabla,\chi_O,\chi_H\}\setminus\{0\}.
\]
If $S$ is nonempty and contains a single sign value, the compatible orientation
is unique:
\begin{equation}
\boxed{
\chi_{\rm TIR}=\chi_\nabla=\chi_O=\chi_H
}
\label{eq:v12-consensus-orientation}
\end{equation}
for every nonzero participating source. States with disagreeing active source
signs are routed to the conflict surface.

\vTwelveStatus{B}{--}{PASS}

\section{What has been forced}
Combining role routing and orientation gives
\begin{equation}
\boxed{
\text{pre-coefficient state}
\longrightarrow
\begin{cases}
R_h\to h,\\
R_a\to a,\\
R_b\to b,\\
R_c\to c,
\end{cases}
\qquad
(\chi_\nabla,\chi_O,\chi_H)
\longrightarrow\chi_{\rm TIR}.
}
\label{eq:v12-coefficient-forcing-chain}
\end{equation}
Thus slot identity and orientation/sign are fixed whenever the declared source
operators are mutually consistent.

\section{Magnitude parent evaluation}
A source audit shows that the historical integration lineage already declares
magnitude-parent expressions for the three charged-lepton action/release
packets. The upstream flavour and Platonic closures supply
\begin{equation}
\boxed{N_F=3,\qquad (L_3,L_4,L_5)=(7,2,5).}
\label{eq:v12-magnitude-upstream-integers}
\end{equation}
Let $I$ denote the unit count, $I=1$, and let $X_i$ denote the typed parent
whose cardinality is $L_i$. Componentwise evaluation of the declared packets
\begin{align}
P_e&=(I,F,I,I),\\
P_{e\mu}&=(0,X_5,X_4,X_3+I),\\
P_{\mu\tau}&=(0,F,I,X_3)
\end{align}
gives
\begin{equation}
\boxed{M(P_e)=(1,3,1,1),}
\label{eq:v12-magnitude-electron}
\end{equation}
\begin{equation}
\boxed{M(P_{e\mu})=(0,5,2,8),}
\label{eq:v12-magnitude-emu}
\end{equation}
and
\begin{equation}
\boxed{M(P_{\mu\tau})=(0,3,1,7).}
\label{eq:v12-magnitude-mutau}
\end{equation}
The evaluation is unique once a parent packet is supplied. It consumes no
measured lepton mass, Yukawa coupling, fitted action or residual-to-target
quantity.

\vTwelveStatus{B}{--}{PASS}
The canonical theorem is
\path{TIR/foundations/TIR_COEFFICIENT_MAGNITUDE_PARENT_EVALUATION_V0_1.md},
and the deterministic validator is
\path{TIR/validation/tir_coefficient_magnitude_parent_evaluation_v0_1.py}.

\section{Residual selector frontier}
The preceding PASS closes arithmetic evaluation of the already-declared parent
expressions; it does not prove that a physical transition uniquely selects those
expressions. The old lineage still marks the transition-level semantic mapping
as \texttt{PROJECT\_MODEL\_ASSIGNMENT}. The remaining coefficient theorem is
therefore the narrower selector
\begin{equation}
\boxed{
\text{coefficient-free transition state}
\longrightarrow
(P_h,P_a,P_b,P_c).
}
\label{eq:v12-coefficient-parent-selector}
\end{equation}
The selector must be evaluated before reading a recovered coefficient tuple or
any mass/Yukawa target.

\vTwelveStatus{B}{--}{OPEN}
This gate is named \texttt{COEFFICIENT\_TRANSITION\_PARENT\_SELECTOR} in
\path{TIR/TIR_COMPLETION_FRONTIER_V0_7.md}.

\section{Circularity firewall}
The coefficient-forcing map admits coefficient-free geometric and routing
parents. Its declared input surface excludes observed particle masses,
observed Yukawa couplings, an already recovered coefficient tuple, a runtime
envelope generated from that tuple, and residual-to-target information. The
retrospective stopping-length pattern remains candidate evidence and is not
promoted by the magnitude parent-evaluation theorem.

\section{Interface to particle sectors}
Historical formulas such as the electron action and the charged-lepton release
operators are migrated to Chapter~\ref{ch:v12_leptons_neutrinos}. There they
are treated as sector constructions over the basis
\eqref{eq:v12-generator-basis}, with their own empirical timing and verdict.
The present chapter owns the upstream generator, typed slot routing, orientation
forcing, exact parent-expression evaluation and the residual transition-parent
selector frontier.

% ---- END TIR/monograph/v12/chapters/ch11_action_architecture_coefficient_forcing.tex ----
\clearpage
